% Bagian: first-order-logic
% Bab: semantics
% Subbagian: intro-semantics

\documentclass[../../../include/open-logic-section]{subfiles}

\begin{document}

\olfileid[id]{fol}{syn}{its}

\olsection{Pendahuluan}

Memberikan makna kepada ungkapan merupakan ranah semantik. Konsep utama
dalam semantik ialah keterpenuhan dalam !!a{structure}. !!^a{structure}
memberikan makna kepada unsur-unsur pembentuk bahasa: !!a{domain} merupakan
himpunan objek yang tak kosong. Kuantor ditafsirkan sebagai merentang atas
  domain ini, !!{constant}s diberi anggota-anggota domain, !!{function}s diberi
  operasi-operasi pada !!{domain}, dan !!{predicate}s diberi
relasi-relasi pada !!{domain}. !!{domain} tersebut beserta penetapan kepada
kosakata dasar membentuk !!a{structure}. !!^{variable}s dapat muncul dalam
!!{formula}s, dan untuk memberikan semantik, kita juga harus menetapkan
!!{element}s dari !!{domain} kepada variabel-variabel itu---inilah penugasan
variabel. Relasi keterpenuhan akhirnya menyatukan semua hal tersebut.
!!^a{formula} dapat terpenuhi dalam !!a{structure}~$\Struct{M}$ relatif
terhadap penugasan !!a{variable}~$s$, yang ditulis sebagai
$\Sat{M}{!A}[s]$. Relasi ini juga didefinisikan melalui induksi pada struktur
dari~$!A$, dengan menggunakan tabel kebenaran bagi penghubung logis untuk
mendefinisikan, misalnya, keterpenuhan $(!A \land !B)$ berdasarkan
keterpenuhan (atau ketakterpenuhan) $!A$ dan~$!B$. Kemudian ternyata bahwa
penugasan !!{variable} tidak relevan jika !!{formula}~$!A$ merupakan
!!a{sentence}, yakni, tidak mempunyai variabel bebas, sehingga kita dapat
berbicara tentang !!{sentence}s yang semata-mata terpenuhi (atau tidak)
dalam !!{structure}s.

Berdasarkan relasi keterpenuhan $\Sat{M}{!A}$ bagi !!{sentence}s, kita
kemudian dapat mendefinisikan gagasan semantis dasar berupa validitas,
perikutan, dan keterpenuhan. !!^a{sentence} valid, $\Entails !A$, jika setiap
!!{structure} memenuhinya. Kalimat itu diakibatkan oleh suatu himpunan
!!{sentence}s, $\Gamma \Entails !A$, jika setiap !!{structure} yang memenuhi
semua !!{sentence}s dalam~$\Gamma$ juga memenuhi~$!A$. Dan suatu himpunan
!!{sentence}s dapat dipenuhi jika ada suatu !!{structure} yang memenuhi
semua !!{sentence}s di dalamnya sekaligus. Karena !!{formula}s didefinisikan
secara induktif, dan keterpenuhan pada gilirannya didefinisikan melalui
induksi pada struktur !!{formula}s, kita dapat menggunakan induksi untuk
membuktikan sifat-sifat semantik kita dan menghubungkan gagasan-gagasan
semantis yang telah didefinisikan.

\end{document}
